\section{Aufmerksamkeit}

\slide{Begrifflichkeiten}{
    \begin{itemize}
        \item \b{Visual Neglect:} Eine Aufmerksamkeitsstörung (keine Wahrnehmungsstörung), bei der die kontraläsionale Seite des Gesichtsfeldes (meist links nach Läsion im rechten Parietallappen) vernachlässigt wird.
        \item \b{Selektiver Filter (Broadbent):} Ein Mechanismus, der Reize basierend auf physikalischen Eigenschaften (z. B. Ort, Stimmlage) selektiert und nur diese zur inhaltlichen Verarbeitung durchlässt ("Alles-oder-Nichts"-Prinzip).
        \item \b{Endogene Steuerung:} Aufmerksamkeit wird aktiv, willentlich und erwartungsgesteuert (Top-Down) ausgerichtet (z. B. durch Interpretation eines Pfeilsymbols).
        \item \b{Exogene Steuerung:} Aufmerksamkeit wird automatisch, reflexartig und reizgesteuert (Bottom-Up) durch saliente Außenreize angezogen (z. B. plötzliches Aufblinken).
        \item \b{Spotlight-Metapher:} Aufmerksamkeit funktioniert wie ein Lichtkegel; Reize innerhalb des Fokus werden verarbeitet, alles außerhalb bleibt im Dunkeln.
    \end{itemize}
}
\slide{Filtertheorien (Broadbent vs. Treisman)}{
    \begin{itemize}
        \item \b{Flaschenhals-Modell (Broadbent):}
        \begin{itemize}
            \item Stufen: Input $\rightarrow$ Sensory Store $\rightarrow$ Selective Filter (Flaschenhals) $\rightarrow$ Higher Level Processing $\rightarrow$ Working Memory.
            \item Prinzip: Frühe Selektion vor der inhaltlichen Analyse; unbeachtete Informationen werden vollständig blockiert.
            \item Evidenz: Split-Span-Paradigma (Schwierigkeit beim schnellen Kanalwechsel).
        \end{itemize}
        \item \b{Dämpfungsmodell (Treisman):}
        \begin{itemize}
            \item Stufen: Input $\rightarrow$ Sensory Store $\rightarrow$ Attenuating Filter (Flaschenhals) $\rightarrow$ Hierarchy of Analyzers $\rightarrow$ Working Memory.
            \item Prinzip: Unbeachtete Informationen werden nicht blockiert, sondern nur gedämpft (abgeschwächt); bei hoher Relevanz (z. B. eigener Name) erreichen sie dennoch die Bewusstseinsschwelle.
            \item Evidenz: Cocktail-Party-Effekt (Wahrnehmung des eigenen Namens im Stimmengewirr) und Shadowing-Experimente (inhaltliche Vermischung bei dichotischem Hören).
        \end{itemize}
        \item \b{Vergleich:}
        \begin{itemize}
            \item Gemeinsamkeit: Beide postulieren einen sensorischen Speicher und einen Engpass (Flaschenhals) aufgrund begrenzter Kapazität.
            \item Unterschied: Broadbent filtert komplett ("Alles-oder-Nichts"), Treisman dämpft nur ("Mehr-oder-Weniger") und erlaubt späte Selektion bei wichtigen Reizen.
        \end{itemize}
    \end{itemize}
}
\slide{Split-Span Paradigma}{
    \begin{itemize}
        \item Ziel: Testet die Wahrnehmungsfähigkeit bei geteilter Aufmerksamkeit (dichotisches Hören) und die Kosten von Kanalwechseln.
        \item Bedingung 1 (Binaural): Gleiche Zahlen auf beiden Ohren $\rightarrow$ keine Selektion nötig, gute Wiedergabe.
        \item Bedingung 2 (Dichotisch, beliebige Wiedergabe): Unterschiedliche Zahlenpaare; Wiedergabe erfolgt meist ohrweise (erst alle links, dann alle rechts) $\rightarrow$ ca. 60\% korrekt (nur ein Filterwechsel nötig).
        \item Bedingung 3 (Dichotisch, paarweise Wiedergabe): Erzwungene Wiedergabe in Präsentationsreih-enfolge (links-rechts, links-rechts) $\rightarrow$ sehr schlecht (20\%), da der Filter nach jeder Zahl umgeschaltet werden muss, was Zeit und Kapazität kostet .
    \end{itemize}
}
\slide{Visuelle Suche (Feature vs. Conjunction)}{
    \begin{itemize}
        \item \b{Feature Search (Merkmals-Suche):}
        \begin{itemize}
            \item Zielreiz unterscheidet sich in einem Merkmal (z. B. Farbe) von Distraktoren.
            \item Prozess: Parallel, automatisch ("Pop-out"-Effekt), reizgesteuert.
            \item Geschwindigkeit: Unabhängig von der Anzahl der Distraktoren (flache Suchkurve).
        \end{itemize}
        \item \b{Conjunction Search (Verbindungs-Suche):}
        \begin{itemize}
            \item Zielreiz ist durch eine \it{Kombination} von Merkmalen definiert (z. B. rotes T unter roten X und schwarzen T).
            \item Prozess: Seriell (Element für Element), kontrolliert, wissensgesteuert; erfordert "Binding" (Zusammenfügen der Merkmale) im Aufmerksamkeits-Spotlight.
            \item Geschwindigkeit: Abhängig von der Anzahl der Distraktoren (steile Suchkurve).
        \end{itemize}
    \end{itemize}
}
\slide{Posner-Paradigma}{
    \begin{itemize}
        \item Aufbau: Reaktion auf einen Zielreiz (Target), dem ein Hinweisreiz (Cue) vorausgeht.
        \item \b{Cues:}
        \begin{itemize}
            \item Endogen (Zentral): Symbolischer Pfeil, erfordert Interpretation und willentliche Ausrichtung.
            \item Exogen (Peripher): Aufleuchten am Zielort, zieht Aufmerksamkeit automatisch an.
        \end{itemize}
        \item Cueing-Effekt: Valide Cues (zeigen korrekten Ort) verkürzen die Reaktionszeit; invalide Cues verlängern sie (Kosten durch Umschalten).
        \item Informativität: Spielt nur bei endogenen Reizen eine Rolle (man folgt dem Pfeil nur, wenn er meistens stimmt); exogene Reize ziehen Aufmerksamkeit auch an, wenn sie oft falsch liegen (Automatismus).
    \end{itemize}
}